\documentclass{article}
\usepackage{graphicx} % Required for inserting images

\usepackage{amsmath}
\usepackage{amssymb}

\newcommand{\fstirling}[2]
    {\left[
    \begin{matrix}
        #1 \\
        #2
    \end{matrix}
    \right]}

\newcommand{\sstirling}[2]
    {\left\{
    \begin{matrix}
        #1 \\
        #2
    \end{matrix}
    \right\}}

\newcommand{\eulerian}[2]
    {\left<
    \begin{matrix}
        #1 \\
        #2
    \end{matrix}
    \right>}

\title{Combinatorics 2}
\author{Siddharth K}
\date{October 2025}

\begin{document}

\maketitle

\section*{Problem 1}
Let $N = p_1p_2\ldots p_n$. Splitting $N$ into its prime factors is the same as splitting the set $\{p_1, \ldots, p_n\}$ into nonempty subsets, which is $B_n = \sum_{k = 0}^n \sstirling{n}{k}$.

\section*{Problem 2}
To create a permutation of $[n]$ with exactly $2$ descents, we can split $[n]$ into three disjoint sets, $A, B, C$, and write each set's elements side by side in ascending order. However, this sometimes creates $1$ or $0$ descents, so we exclude those from our total count. \\

The number of ways to split $[n]$ into three subsets is $3^n$. However, if every element in $A$ is less than every element in $B$, whose elements are less than every element in $C$, the permutation does not have $2$ descents. The number of these permutations is $(n + 1)^2$, since we choose two points between $0$ and $n$ to be the splitting point between $A$, $B$ and $B$, $C$. \\

Also, if every element in $A$ is less than every element in $B$, the permutation doesn't have $2$ descents. We can count the number of these sets by fixing the size of $C$ to be $k$. Then there are $n - k + 1$ ways to choose the split point between $A$ and $B$ and $\binom nk$ ways to choose $C$, so the total count is
$$\sum_{k = 0}^n \binom nk (n - k + 1) = \sum_{k = 0}^n \binom nk n - \sum_{k = 0}^n \binom nk k + \sum_{k = 0}^n \binom nk$$
$$= (n + 1)\sum_{k = 0}^n \binom nk - n\cdot 2^{n - 1} = (n + 1)\cdot 2^n - n\cdot 2^{n - 1}$$
$$= (n + 2)\cdot 2^{n - 1}.$$
This is also the number of ways to make the split where every element in $B$ is less than every element in $C$. However, summing both cases overcounts the case where all of $A$ is less than all of $B$ which is also less than all of $C$, so we subtract that count once. \\

The overall sum is
$$3^n - (2\cdot (n + 2)\cdot 2^{n - 1} - (n + 1)^2)$$
$$= 3^n - 2^n(n + 2) + (n + 1)^2.$$

\section*{Problem 3}
For the lower bound, we can make a set partition of $[2n]$ into $[n]$ blocks as follows. First, start with the blocks $\{1\}, \{2\}, \ldots, \{n\}$. Then, for each element $n + 1, \ldots, 2n$, choose one of the $n$ blocks to drop it into. Then the resulting partition contains $n$ nonempty blocks, so it is counted by $\sstirling{2n}{n}$. The number of ways to do this is $n^n$, so for $n \ge 2$,
$$\sstirling{2n}{n} \ge n^n > n!.$$
For the upper bound, we form a permutation of $[2n]$ given a set partition of $[2n]$ into $n$ blocks. First, sort each block internally in increasing order, and then order all the blocks by their minimum element. Finally, concatenate the blocks to form a permutation of $[2n]$. Each set partition produces a unique permutation, so $\sstirling{2n}{n} \le (2n)!$. Furthermore, the inequality is strict for $n \ge 2$, since a permutation starting with $2, 1, 3, 4, \ldots$ can't be produced in this way.

\section*{Problem 4}
The number of set partitions of $[n]$ with $\{1\}$ as one of the subsets is $B_{n - 1}$, so the number of set partitions of $[n]$ without $\{1\}$ is $B_n - B_{n - 1}$.

\section*{Problem 5}
First, we count the number of ways to split $[n]$ into two parts, which is $\sstirling n2$, and subtract the number of ways to split $[n]$ into two subsets of equal size, which is $\binom{n}{n/2}$. So the result is $\sstirling n2 - \binom{n}{n/2}$.

\section*{Problem 6}
Let $a_n$ be the number of permutations that have $1$ and $2$ in the same cycle. Given a permutation of $[n - 1]$ with $1$ and $2$ in the same cycle, we can form a permutation of $[n]$ by placing $n$ in front of any of the existing numbers, or we can place it in its own cycle. So $a_n = (n - 1)a_{n - 1} + a_{n - 1} = na_{n - 1}$ for all $n \ge 3$. We have that $a_1 = 0$ and $a_2 = 1$, so $a_3 = 3$, $a_4 = 12$, and so on. In general, $a_n = n!/2$. \\

Let $b_n$ be the number of permutations that have $1$, $2$, and $3$ in the same cycle. Similar to before, $b_n = nb_{n-1}$ for $n \ge 4$. We have that $b_1 = b_2 = 0$ and $b_3 = 3!/3 = 2$, so $b_4 = 8$, $b_5 = 40$, and in general, $b_n = n!/3$. \\

Let $c_n$ be the number of permutations with $1$ and $2$ in the same cycle and $3$ and $4$ in the same cycle. We have that $c_n = nc_{n - 1}$ for all $n \ge 5$. Also, $c_1 = c_2 = c_3 = 0$, and $c_4 = 1 + 4!/4 = 7$. So $c_n = n(n - 1)\cdots(6)(5) \cdot 7 = n! \cdot \frac{7}{24}$.

\section*{Problem 7}
Assume the $n$ people sit at $i$ tables. Then the number of ways to seat them is $\fstirling ni$ and the number of ways for all of them to choose a menu item is $k^i$. Summing over $i$, we get a total of
$$\sum_{i = 0}^k \fstirling ni k^i.$$

\section*{Problem 8}
Take any set partition of $[n]$ (this can be done in $B_n$ ways). If the partition has no blocks of size $1$, keep it as is (counted by $a_n$). Otherwise, merge all the blocks of size $1$ into one block and add $n + 1$ (counted by $a_{n + 1}$). This process is reversible, since we can take the block with $n + 1$ and split its elements into individual blocks. So we have a bijection, and $B_n = a_n + a_{n + 1}$.

\section*{Problem 15}
\subsection*{Part (a)}
The LHS is the number of ways to split $[n + 1]$ into $m + 1$ disjoint nonempty subsets. We can also count this by the size of the block that contains $n + 1$. We can choose $k$ elements that are not in that block in $\binom nk$ ways and partition the remaining k elements in $\sstirling km$ ways. Summing over $k$ gives us the RHS.

\subsection*{Part (b)}
The LHS is the number of permutations of $[n + 1]$ with $m + 1$ cycles. To count this, we can start with a permutation of $[n]$ with $k$ cycles, and choose $m$ of them to be left as separate cycles. We can concatenate the other $k - m$ cycles into one big cycle along with the new element $n + 1$. This is reversible, since we can delete $n + 1$ and split its cycle into the $k$ original cycles. So the LHS and RHS are equal.

\section*{Problem 17}
Each card in a deck can be mapped to a set partition of $[5]$. I was thinking of something where we ask the spectator to secretly split five different objects into a few groups, then we ask a few identifying questions (as seemingly unrelated as possible) to isolate a card, and then map that card to the spectator's set partition. However I couldn't come up with anything more specific.

\section*{Problem 18}
\subsection*{Part (a)}
Assume $f(x) = x^m$. Simplifying the LHS, we get
$$x\frac{d}{dx}(x^m) = x(mx^{m -1}) = mx^m,$$
$$\left(x\frac{d}{dx}\right)^n(x^m) = m^nx^m.$$
Simplifying the RHS,
$$x^k\frac{d^k}{dx^k}(x^m) = m(m - 1)\cdots (m-k + 1)x^m = m^{\underline{k}}x^m,$$
$$\sum_{k = 0}^n \sstirling nk x^k \frac{d^k}{dx^k}(x^m) = \sum_{k = 0}^n \sstirling nk m^{\underline{k}}x^m = m^nx^m.$$
So the LHS and RHS are equal. Since summations and the derivative operator are linear, this applies to all polynomial functions $f$.

\subsection*{Part (b)}
Again, assume $f(x) = x^m$. The LHS becomes
$$x^n\frac{d^n}{dx^n}(x^m) = m^{\underline{n}}x^m.$$
The RHS becomes
$$\sum_{k = 0}^n (-1)^{n - k} \fstirling nk x^k\frac{d^k}{dx^k}(x^m) = \sum_{k = 0}^n (-1)^{n - k} \fstirling nk m^kx^m = m^{\underline{n}}x^m.$$
So the LHS and RHS are equal for all polynomial functions $f$.

\end{document}